%!TEX program = xelatex
\documentclass[10pt,twocolumn]{article}

\usepackage{amsmath,amsfonts,amssymb,amsthm,mathtools,stmaryrd}
\usepackage{booktabs,fancyhdr}
\usepackage[shortlabels]{enumitem}
\usepackage[%
%margin=1in,
top=0.65in,bottom=0.75in,left=0.5in,right=0.5in,landscape
]{geometry}
\setlength{\columnsep}{0.5in}


\usepackage[no-math]{fontspec}
\setmainfont{TeX Gyre Pagella}

\usepackage{pifont}
\newcommand{\cmark}{\ding{51}}% new checkmark
\newcommand{\xmark}{\ding{55}}% new x-mark

\newtheorem*{theorem}{Theorem}
\theoremstyle{definition}
\newtheorem*{definition}{Definition}

\usepackage[%
natbib=true,backend=biber,sorting=nymdt,%
maxbibnames=10,%
citestyle=authoryear,bibstyle=authoryear,%
url=false,doi=false,isbn=false%
]{biblatex}
\input{tools/jbib_links.tex}
\addbibresource{test.bib}
\input{tools/jformat.tex} % note: imports hyperref
\usepackage{cleveref}

% warning character
\usepackage{newunicodechar}
\newcommand{\Warning}{%
  \makebox[1.4em][c]{%
    \makebox[0pt][c]{\raisebox{.15em}{\small!}}%
\makebox[0pt][c]{\color{red}\Large$\triangle$}}}%

\newcommand\hlbox[2][]{%
  \tikz[overlay]\node[%
  fill=yellow!20,inner sep=0pt,
  anchor=text,rectangle,rounded corners=1mm,#1]%
{#2\strut};\phantom{#2}}
\newcommand\hlA[1]{\hlbox[fill=blue!10]{#1}}
\newcommand\hlB[1]{\hlbox[fill=red!10]{#1}}
\usepackage{soul}

\usepackage[linguistics]{forest}
% For putting parens around a subtree
\usetikzlibrary{calc}
\tikzset{
  ncbar angle/.initial=90,
  ncbar/.style={%
    to path= (\tikztostart)
    -- ($ (\tikztostart)!#1!\pgfkeysvalueof{/tikz/ncbar angle}: (\tikztotarget)$)
    -- ($ (\tikztotarget)! ($(\tikztostart)!#1!\pgfkeysvalueof{/tikz/ncbar angle}: (\tikztotarget)$)!\pgfkeysvalueof{/tikz/ncbar angle}: (\tikztostart)$)
    -- (\tikztotarget)
  },
  ncbar/.default=0.5cm,
}
\tikzset{left brace/.style={ncbar=-0.5mm}}
\tikzset{right brace/.style={ncbar=0.5mm}}
\tikzset{left paren/.style={ncbar=0.5mm,out=-110,in=110}}
\tikzset{right paren/.style={ncbar=0.5mm,out=-70,in=70}}

\input{tools/forest-trees.tex}
\input{tools/kast.tex} % for a nice box
\input{tools/gb4e-examples.tex} % import gb4e, tweak
\setlength{\intextsep}{0ex}

%%%%  Author name
\def\FNAME{Jacob L}
\def\LNAME{Hover}

\pagestyle{fancy}
\fancyhead[L]{\thetitle}
\fancyhead[R]{\LNAME}

\begin{document}
\title{Linguistics handout template}
\author{\FNAME\ \LNAME}
\date{Spring 2021}

\maketitle

%%%%%%%%%%%%%%%%%%%%%%%%%%%%%%%%%%%%%%%%%%%%%%%%%

\section{Introduction}%
\label{sec:introduction}

Lorem ipsum here's a citation \citep{bombar}. \Citet{nymous} says something. \textbf{Note automatic hyperlinks in the references.}
Use \texttt{mdframed} for a simple box. 

\begin{mdframed}[style=rounded]
Here's a description:
  \begin{description}
    \item[Generalization] Yup.
    \item[Reported exceptions] A \hlbox{highlighted} thing.
  \end{description}
\end{mdframed}

\begin{exe}
  \ex\tree[for tree={l sep=1.5ex, l=1.5ex, inner sep=0}]{%
    [\textit{v}P 
    [{\textit{v}\\\textit{-ya}\\([\textsc{instr}])}] 
  [VP [V] [DP]]]}
\end{exe}

The \texttt{kast} environment gives a box with a title.

\begin{kast}{Proposal for X}

  \emph{kast} is Estonian for `box'.

  Include a graphic

  \noindent\includegraphics[width=.95\textwidth]{example.png}
\end{kast}

% NOTE: If you want to put a graphic inside an itemize environment like this, you may want to `clip' it like this:
% \begin{itemize}
%   \item Here's an item
%   %\includegraphics[width=.85\textwidth/2,trim={15mm 0 0 0},clip]{example.png}
% \end{itemize}


\subsection{A subsection}
Here's a tree with some Ti\textit{k}Z annotation.
\begin{exe}
  \ex Another tree\\
  \tree{%
    [Voice$_\textsc{pass}$P, for tree={s sep=15pt}
    [Voice$_\textsc{pass}$P
    [Voice$^0_\textsc{pass}$\\$\theta$]
    [VP, for tree={s sep=15pt}]]
    [PP,edge=dashed, tikz={
      \node(x) [inner sep=0pt,fit to=tree] {};
      \draw[] (x.north west) to [left brace] (x.south west);
      \draw[] (x.north east) to [right brace] (x.south east);
    }
  [P] [DP\\\textit{agent}]]]}
\end{exe}
%%%%%%%%%%%%%% REFERENCES %%%%%%%%%%%%%%%%%%%%%%%

\defbibnote{note:complete}{\it See bibliography in \cite{nymous} for further references.}
%\AtNextBibliography{\footnotesize}% For resizing the bibliography
\printbibliography[%
title=Selected References,%
prenote=note:complete% define a bibnote to print before the bibliography
]{}

%\noindent\hrulefill%

%%%%%%%%%%%%%  APPENDIX   %%%%%%%%%%%%%%%%%%%%%%%%
\appendix

\section{Afterthoughts}

\begin{kast}{Good questions}
  \begin{description}
    \item[Jiminy asks]  This is a question?
  \end{description}
\end{kast}

\end{document}